\documentclass[11pt,a4paper]{article}
\usepackage[utf8]{inputenc}
\usepackage{amsmath}
\usepackage{amssymb}
\usepackage{amsfonts}
\usepackage{graphicx}
\usepackage{hyperref}
\usepackage{geometry}
\geometry{margin=1in}

\title{tagd Mathematical Foundations}
\author{Isaac N Colley}
\date{April 2026}

\begin{document}

\maketitle

\begin{abstract}
This document formally describes the mathematics already present in the tagd semantic-relational system and the TAGL language. It builds strictly from the source code, using only concepts observable in the implementation. No new mathematics is invented.
\end{abstract}

\section{Introduction}

A \textbf{tagspace} is a finite set of tags together with two kinds of relations:

\begin{itemize}
    \item \textbf{Identity relations} (vertical): exactly one subordinate relation per tag, inducing a rooted tree via ranks.
    \item \textbf{Predicate relations} (horizontal): zero or more additional labeled relations attached to each tag.
\end{itemize}

All semantics derive from structure. The system is deterministic and inspectable because its canonical form (\texttt{dump()}) is a pure textual representation of this structure.

\section{Primitives}

\subsection{Tag}
A tag consists of:
\begin{itemize}
    \item a globally unique UTF-8 label \texttt{\_id} (nominal identity)
    \item exactly one subordinate relation (structural identity)
\end{itemize}

\subsection{Subordinate Relation (Identity Relation)}
A subordinate relation is the pair:
\[
(sub\_relator,\ super\_object)
\]
where both are tags.

\textbf{Invariant} (enforced by source): Every tag except the root \texttt{\_entity} has exactly one such relation.

\subsection{Rank}
The \texttt{rank} of a tag is a sequence of UTF-8 code points.

\textbf{Core rule} (from \texttt{rank.h} and \texttt{rank::contains()}):
\[
A\ \_sub\ B \iff \operatorname{rank}(B)\ \text{is a prefix of}\ \operatorname{rank}(A)
\]

This defines \textbf{containership}: every descendant’s rank contains its ancestor’s rank as a prefix.

\section{Core Model -- The Tagspace as a Poset}

A tagspace \( (T, \preceq) \) is a \textbf{partially ordered set (poset)} under the containership order:
\[
A \preceq B \iff \operatorname{rank}(B)\ \text{is a prefix of}\ \operatorname{rank}(A)
\]

\begin{itemize}
    \item The order is reflexive, antisymmetric, and transitive.
    \item It forms a \textbf{rooted tree} with root \texttt{\_entity} (the unique minimal element, rank \( \epsilon \)).
\end{itemize}

Source reference: \texttt{rank.h}, \texttt{abstract\_tag::operator<}, and \texttt{merge\_containing\_tags()} in \texttt{tagd.cc}.

\section{Topology Induced by the Poset (Alexandrov Topology)}

The poset induces a natural topology called the \textbf{Alexandrov topology}:

\begin{itemize}
    \item \textbf{Open sets} = upward-closed sets = subtrees (if a tag is included, all its descendants are included).
    \item \textbf{Closed sets} = downward-closed sets = paths from the root to a tag.
\end{itemize}

This matches the source exactly: queries and \texttt{rank::contains()} operate on subtrees. The topology encodes reachability and inheritance directly from the rank prefix order.

\section{Hard Tags as Universal Base}

Every tagspace shares the same fixed hard tag subtree (from \texttt{hard-tags.h} and bootstrap). This is the \textbf{invariant base}:

\begin{itemize}
    \item \texttt{\_entity} is the root.
    \item Hard tags (\texttt{\_sub}, \texttt{\_rel}, \texttt{\_number}, \texttt{\_interrogator}, etc.) form the common skeleton.
    \item Every user-defined tag extends this base.
\end{itemize}

Mathematically, the hard tag structure is a \textbf{fixed subspace} that every tagspace contains by construction. Any map between tagspaces must preserve this base.

\section{Predicates -- Horizontal Structure}

Predicate relations (\texttt{\_has}, \texttt{\_can}, \texttt{\_refers\_to}, etc.) are additional labeled relations attached to tags. They do \textbf{not} affect the poset order or rank.

They sit on top of the topological skeleton as extra structure on the nodes.

\section{Lexical Structure and Syntax}

TAGL is the language for constructing and querying the poset + predicates:

\begin{itemize}
    \item \texttt{>> subject sub\_relator super\_object} adds a new element to the poset.
    \item Predicates attach horizontal relations without changing position in the tree.
\end{itemize}

The grammar in \texttt{parser.y} enforces the single-identity invariant at parse time.

\section{Query Semantics}

Queries traverse the poset using the containership order:
\begin{itemize}
    \item Subtree queries (\texttt{\_sub}) return upward-closed sets.
    \item Wildcard \texttt{*} matches any relator within the subtree.
\end{itemize}

This is exactly the topological notion of open sets generated by the partial order.

\section{Canonical Form and Serialization}

The \texttt{dump()} method produces a total ordering of the poset that respects the rank prefix order, yielding a canonical textual representation of the entire structure.

\section{Mathematical Summary}

A tagspace is a structure \( (T, \preceq, P) \) where:

\begin{itemize}
    \item \( (T, \preceq) \) is a finite rooted tree poset equipped with the Alexandrov topology induced by rank prefixes.
    \item \( P \) is a set of labeled predicate relations on \( T \).
    \item Every element has exactly one covering relation in \( \preceq \).
    \item The hard tag subtree is the universal invariant base shared by all tagspaces.
\end{itemize}

TAGL is the language for constructing and querying this structure.

\section{References (Source of Truth)}

\begin{itemize}
    \item \texttt{rank.h} — rank prefix and containership
    \item \texttt{tagd.cc} — poset construction, \texttt{merge\_containing\_tags}
    \item \texttt{parser.y} — grammar enforcing single identity
    \item \texttt{TAGL-spec.md} — language definition
    \item \texttt{hard-tags.h} — axiomatic universal base
\end{itemize}

This document will be extended section-by-section as further mathematical insights are extracted from the source. All statements are directly traceable to the code.

\end{document}

